\chapter{Discussion}
\label{ch:discussion}

This chapter confronts the objections a competent examiner should raise ---
each stated in its strongest form, answered with committed evidence, and
never softened. The full pre-answered set is maintained in
\evfile{docs/DEFENSE\_QA.md}; the load-bearing ones follow.

\section{Threats to validity}
\label{sec:disc-validity}

\paragraph{``Your headline numbers come from your own simulator.''}
True, and framed that way everywhere: decision quality under a stated system
model. Three things bound the gap. First, \emph{calibration direction}: the
congestion form was fitted against a real inference server and the model
over-penalizes high utilization ($a{=}0.86$ vs.\ assumed $1.0$) --- it errs
against \polyforge{}'s lean postures; the GPU tier bench confirms tier
ordering on real hardware. Second, \emph{real data}: the cost and
forecasting claims are restated on replayed BurstGPT demand ($n{=}96$
confirmatory) and the cache claims are measured on real LMSYS prompts.
Third, the \emph{live ordinal check}: the kind-cluster experiment asks only
whether the simulator's ranking reproduces, with the HPA arm verified and
the \jcac{} arm gated behind an executable actuation check. What is not
claimed: absolute dollars or latencies transferring to any production
fleet.

\paragraph{``You designed the v3 overload cells so your system would
win.''} The cells were sized by arithmetic on committed model constants,
pinned by unit tests, not tuned by trial runs. Three protections: the
pre-registration named the sole treatment and hypotheses before any run; a
same-envelope control cell (\texttt{ramp\_gentle}) exists where the alleged
advantage should not appear --- and it does not; and the confirmatory H1$'$
still \emph{failed}. A designed regime with an honest null inside it is
evidence of measurement, not of rigging.

\paragraph{``$p=10^{-33}$? Really?''} At hundreds of paired seeded simulator
runs, tiny p-values measure the simulator's determinism, not the effect. The
citation unit is the paired effect size with a bootstrap CI
(Section~\ref{sec:method-stats}); p-values appear once, as prereg gate
outcomes.

\paragraph{``Your SLO story moved the goalposts.''} The ledger says
otherwise, and the ledger is the answer: seven pre-registrations pushed
before their runs, two confirmatory SLO failures published as failures, one
mechanism confirmation, and a stopping rule forbidding a third attempt. The
final SLO claim --- parity at $-43\ldots{-48\%}$ cost --- is what survived,
not what was hoped for. Goalpost-moving is invisible when it happens in
private; here it happened in public, in the git history, where it can be
diffed.

\paragraph{``Were the baselines actually tuned, or strawmen?''}
Grid-searched on the paper's own composite objective, per baseline, sweeps
committed. The grids are bounded at a vendor-sane envelope because $J$
improves monotonically toward over-provisioning --- an unbounded grid would
tune every baseline into a cost blowout and \emph{flatter} \polyforge{}.
Iso-cost variants remove the ``it just spends less'' objection, and their
gate is reported FAIL-honest where it fails.

\section{Construct scope}

\paragraph{Fairness altitude.}
\label{sec:disc-fairness-altitude}
Jain's index over SLO satisfaction is not fairness as OSDI defines it: VTC
and successors define \emph{service} fairness at token granularity inside a
continuous-batching engine. \polyforge{}'s fairness is capacity allocation
across tenants at the control plane, the claims are scoped to that altitude,
and the \texttt{vtc\_replica} baseline is a replica-level transplant of
VTC's least-service-first idea, stated as such. Beating it on Jain says
\polyforge{} allocates capacity more evenly at lower cost --- not that it
schedules tokens more fairly. Token-level fair schedulers would run
\emph{inside} each replica pool \polyforge{} sizes; the two are
complementary.

\paragraph{Two economies in one cost model.} Replica-hours are priced as
infrastructure (\$0.048/replica-hr; the live arm meters real
replica-seconds) while model-tier calls are priced at market API ratios
(1:10:100). A self-hosting operator should read the tier ratios as relative
GPU-time costs --- the tier bench measured 1:1.5:16.6 latency ratios on real
hardware: same ordering, compressed scale, disclosed. Spot pricing, MIG
partitioning, and fractional GPUs are below the model's altitude and named
out of scope. Reconfiguration friction is not ignored: the realism ablation
bills startup lag and cold caches; multi-minute large-model pulls exceed
what it models, and the limitation says so.

\paragraph{Cache hits that are wrong answers.} Anticipated and measured
(Section~\ref{sec:eval-cache}): precision with strata and a
per-1k-incorrect-hits figure next to the savings, plus the stochasticity
ceiling that bounds what the proxy can say. No ``quality-adjusted dollar''
is invented, because pricing a wrong answer is application-specific.
Staleness and TTL are out of scope and listed as such.

\section{Limitations}
\label{sec:disc-limits}

\begin{itemize}
  \item \textbf{Substrate.} Headline numbers are simulator decision quality;
        live validation is ordinal-only, currently verified for the HPA arm
        with the \jcac{} live slice pending
        (Section~\ref{sec:eval-live}).
  \item \textbf{The live figure exercises one knob.} The live data plane
        burns fixed CPU per request kind, so cache and tier are inert there;
        the live check validates the replica-control projection, with the
        other knobs' realism carried by their own real-data protocols. The
        frozen caption states this.
  \item \textbf{Latency is p95, request-level.} The simulator's tail beyond
        p95 is not credible and is not reported; token-level dynamics (TTFT,
        decode-length tails, head-of-line blocking under continuous
        batching) are outside the model. The live path can report any
        percentile; p99 live restatement is future work, not silently
        substituted.
  \item \textbf{Trace scope.} Demand claims rest on BurstGPT plus the Azure
        LLM 2024 release; prompt claims on LMSYS-Chat-1M. One
        generalization failure is already published (the seasonal
        non-transfer), which is the strongest available evidence that the
        remaining claims are scoped honestly. FaaS/VM traces were
        deliberately descoped as non-LLM demand.
  \item \textbf{Tenant count.} Eight tenants per cluster is the studied
        composition factor; the planner's measured scaling (exponent 1.45,
        timeout crossover at 128 tenants) bounds, but does not remove, the
        gap to production tenant populations. Coordination costs outside
        the planner are unmeasured.
  \item \textbf{Production scale.} SageServe-class fleets serve more in a
        day than the entire replay; the 10.63M-request demand replay is the
        honest analog at decision level, and no stronger claim is made.
\end{itemize}

\section{The honest-nulls ledger as a contribution}

Seven published negatives travel with the wins: v2 H1, v3 H1$'$, the
underpowered first replay, the $\gamma$-term null, the seasonal
non-transfer, the structurally impossible raw per-metric gate, and the cache
hit-quality result. None was hidden; several redirected the design (Holt as
recommended default; per-stream forecasting; per-tenant caches). The
methodological claim of this thesis --- pre-registration with published
nulls in a systems evaluation --- is independent of any performance number
and, in the surveyed literature, unique to this work.
